\input{../tex-inputs/latex-leseansicht-vorspann}

\section[Arthur, Olga und Louise Schnitzler an Gustav Schwarzkopf, 9. 2. 1908]{L04391 Arthur, Olga und Louise Schnitzler an Gustav Schwarzkopf, 9. 2. 1908}
\nopagebreak\mylabel{L04391v}
\rehead{ }\normalsize\beginnumbering\briefempfaengerindex{Schwarzkopf, Gustav@\textsc{Schwarzkopf, Gustav}!zzzSchnitzler, Louise@\emph{von Louise Schnitzler}!1908-02-091@{9. 2. 1908}|(be}\briefempfaengerindex{Schwarzkopf, Gustav@\textsc{Schwarzkopf, Gustav}!zzzSchnitzler, Olga@\emph{von Olga Schnitzler}!1908-02-091@{9. 2. 1908}|(be}\briefempfaengerindex{Schwarzkopf, Gustav@\textsc{Schwarzkopf, Gustav}!zzzSchnitzler, Arthur@\emph{von Arthur Schnitzler}!1908-02-091@{9. 2. 1908}|(be}
\toendnotes[C]{\smallbreak\pagebreak[2]}
\correspDesc{Versand  durch Arthur Schnitzler, Olga Schnitzler, Louise Schnitzler am 9. 2. 1908 in Semmering
\newline{}Übermittlung  am 10. 02. 1908 in Semmering
\newline{}Erhalt  durch Gustav Schwarzkopf im Zeitraum [10. 2. 1908 – 14. 2. 1908?] in Wien}\toendnotes[C]{\smallbreak}
\Standort{DLA, A:Schnitzler, HS.1985.1.1897.}
\physDesc{Bildpostkarte, 288 Zeichen
\newline{}Handschrift Arthur Schnitzler: Bleistift, deutsche Kurrent
\newline{}Handschrift Olga Schnitzler: Bleistift, lateinische Kurrent
\newline{}Handschrift Louise Schnitzler: Bleistift, deutsche Kurrent
\newline{}Versand: Stempel: »\nobreak{}\oindex{Semmering@\textbf{Semmering}|pwk}Semm{[}eri{]}n{[}g{]}, 10. II. 08, XII\nobreak{}«.  }\toendnotes[C]{\smallbreak}\pstart{}{\pb}Hrn \textsc{Gustav Schwarzkopf}\pend{}\pstart{}\textsc{Wien} I\oindex{Wien@\textbf{Wien}|pw}\pend{}\pstart{}\textsc{Tiefer Graben 23}\oindex{Wien@\textbf{Wien}!I., Innere Stadt@\textbf{I., Innere Stadt}!Tiefer Graben 23@\textbf{Tiefer Graben 23}|pw}\pend{}{\bigskip}
\pstart
           \centering{}{\pb}\textcolor{gray}{\textbf{Semmering\oindex{Semmering@\textbf{Semmering}|pw}. Südbahnhotel\oindex{Südbahnhotel [Semmering]@\textbf{Südbahnhotel [Semmering]}|pw}.}}\pend
           \vspace{1em}
\pstart
           {\pb}9/2 908\pend
           
\pstart{}Lieber Guſtav,\pend\vspace{0.5em}
\pstart
           kommen Sie doch! Wir bleiben noch über 8 Tage. Es iſt ſehr ſchön.\pend
           
\pstart
           herzlichſt Ihr{\\[\baselineskip]}\spacefill\mbox{Arthur}\pend
           \leftskip=0em{}\selectlanguage{ngerman}\vspace{1em}
\pstart
           \noindent{}{[}hs. O. Schnitzler:{]} \label{T_L04391-1v}\edtext{Es wäre sehr schön, Sie hier zu sehen,}{\lemma{\textnormal{\emph{Es … sehen,}}}\Cendnote{\textnormal{Olga Schnitzlers\pwindex{Schnitzler, Olga 17.\,1.\,1882 Wien – 13.\,1.\,1970 Lugano@\textsc{Schnitzler, Olga} (17.\,1.\,1882 Wien – 13.\,1.\,1970 Lugano)|pwk} Gruß ist am rechten Rand des Textfeldes um 90° gedreht platziert.}}}\label{T_L04391-1} und es würde
               Ihnen gewiss sehr wol tun.\pend
           
\pstart
           Herzliche Grüsse.{\\[\baselineskip]}\spacefill\mbox{OlgaS.}\pend
           \leftskip=0em{}\selectlanguage{ngerman}\vspace{1em}
\pstart
           \noindent{}{\pb}{[}hs. L. Schnitzler:{]} Herzliche Grüße{ }ſendet\pend
           \pstart \spacefill\mbox{Louise Schnitzler}\pend{}\selectlanguage{ngerman}\endnumbering\briefempfaengerindex{Schwarzkopf, Gustav@\textsc{Schwarzkopf, Gustav}!zzzSchnitzler, Louise@\emph{von Louise Schnitzler}!1908-02-091@{9. 2. 1908}|)be}\briefempfaengerindex{Schwarzkopf, Gustav@\textsc{Schwarzkopf, Gustav}!zzzSchnitzler, Olga@\emph{von Olga Schnitzler}!1908-02-091@{9. 2. 1908}|)be}\briefempfaengerindex{Schwarzkopf, Gustav@\textsc{Schwarzkopf, Gustav}!zzzSchnitzler, Arthur@\emph{von Arthur Schnitzler}!1908-02-091@{9. 2. 1908}|)be}\mylabel{L04391h}
\begin{anhang}
\end{anhang}\newcommand{\dateiname}{L04391}\newcommand{\titel}{Arthur, Olga und Louise Schnitzler an Gustav Schwarzkopf, 9. 2. 1908}\newcommand{\editorInnen}{Herausgegeben von Jahnke, SelmaMüller, Martin Anton}\input{../tex-inputs/latex-leseansicht-abspann}
